\documentclass[11pt,a4paper]{article}
\usepackage[utf8]{inputenc}
\usepackage{ctex} % 支持中文，在 Overleaf 中选择 XeLaTeX 或 LuaLaTeX 即可完美编译
\usepackage{amsmath}
\usepackage{amssymb}
\usepackage{amsfonts}
\usepackage{physics} % 提供优雅的数学物理公式符号（如 \grad, \laplacian, \pdv 等）
\usepackage{geometry}
\geometry{left=2.5cm,right=2.5cm,top=3cm,bottom=3cm}
\usepackage{booktabs}
\usepackage{graphicx}
\usepackage{hyperref}
\usepackage{xcolor}
\usepackage{titlesec}
\usepackage{fancyhdr}
\usepackage[ruled,linesnumbered,vlined]{algorithm2e} % 优雅的伪代码排版

% 配色方案与超链接设置
\definecolor{navyblue}{RGB}{20, 50, 110}
\hypersetup{
    colorlinks=true,
    linkcolor=navyblue,
    citecolor=navyblue,
    urlcolor=navyblue
}

% 算法关键字汉化
\SetAlgorithmName{算法}{算法}{算法列表}
\SetKwInput{Input}{输入}
\SetKwInput{Output}{输出}
\SetKw{To}{to}

% 页眉页脚设置
\pagestyle{fancy}
\fancyhf{}
\rhead{\small 基于 SPH 与顺推冲量的高性能 2D 物理沙盒系统}
\lhead{\small 学术与工程报告}
\cfoot{\thepage}

\title{\vspace{-1.5cm}\textbf{\Large 基于 SPH 流体与顺序冲量凸刚体解算的高性能 2D 物理沙盒系统}\\ \large 稳定性深度优化与底层计算架构重构研究报告}
\author{\textbf{刘进杰} \quad \textbf{侣云飞} \\ \small 计算机科学与技术基地班}
\date{\today}

\begin{document}

\maketitle

\begin{abstract}
本文针对同屏上万个流体粒子与多边形凸刚体进行强交互仿真时极易出现的穿模（Tunnelling）、高压堆积数值爆炸（Blowup）以及多核并发下帧率断崖式退化等物理稳定性与底层计算瓶颈进行了系统性的研究与架构重构。在物理算法层面，本系统采用了弱可压缩平滑粒子流体动力学（WCSPH）与二阶显式速度 Verlet 积分器，并融合了基于符号距离场（SDF）与连续碰撞检测（CCD）的组合式硬边界，同时在冲量累积层面设计了速度锁截断（Velocity Clamping），彻底杜绝了高压与高速下的物体穿模逃逸。在底层计算与渲染架构层面，本文抛弃了传统的 CPU 栅格化即时指令渲染管线，重构为基于点精灵（Point Sprites）硬件加速的 QOpenGLWidget 实例化渲染机制，并且设计了基于 Counting Sort 的 $O(N)$ 空间哈希重建算法、Cache Line 友好的物理坐标 Entry 预取胖节点以及 OpenMP 无竞争堆内存分配结构。实验数据证明，重构后的系统在同屏释放 $5000+$ 个流体粒子及大量凸刚体物体的极端场景下，依然能够保持 $38+$ FPS 的极高交互频率，展现出极强的数值鲁棒性与计算吞吐能力。

\textbf{关键词：} 平滑粒子流体动力学；凸刚体动力学；顺序冲量法；速度 Verlet 积分；连续碰撞检测；缓存行预取；点精灵渲染
\end{abstract}

\noindent\rule{\textwidth}{0.5pt}

\section{引言}
实时多物理场耦合仿真系统（包含流体与凸刚体动力学）是现代计算机图形学与游戏物理引擎的最前沿课题之一。流体通常采用拉格朗日框架下的无网格平滑粒子流体动力学（SPH）模型，其具备描述流体飞溅、破碎和自由表面流动的天然优势；刚体则通过牛顿-欧拉方程与单侧非穿透接触约束限制相对运动，利用顺序冲量法（Sequential Impulses）在每个物理子步迭代求解冲量反射。

然而，当仿真场景向密集交互（即大量粒子高压堆积在刚体上、物体高速碰撞边界）以及大规模实时渲染拓展时，传统的数值积分与碰撞响应系统在有限的时间步长下往往会发生灾难性的崩溃：
\begin{enumerate}
    \item \textbf{高速穿模（Tunnelling 隧道效应）}：超出了离散碰撞检测（DCD）的安全时间窗，物体在一帧之内瞬间越过墙壁边界；
    \item \textbf{堆积数值爆炸（Instability Blowup）}：在高密度堆积环境下，SPH 的压力力学与刚体冲量累加会呈非物理的指数级增长，速度发散导致整个世界完全脱靶；
    \item \textbf{帧率骤降与计算瓶颈}：大量的动态内存碎片分配、多线程锁冲突以及 CPU 栅格化渲染数万粒子的指令耗时，拖垮了整个程序的实时交互表现。
\end{enumerate}

针对上述稳定性问题与计算天花板，刘进杰与侣云飞在计算机科学与技术基地班的研究框架下，对本 2D 高性能物理沙盒系统进行了底层的彻底升级。本文将详尽推导其背后的数学物理方程式，论述穿模、爆炸和帧率降低的根因，并给出相对应的重构方案。

\section{物理系统数学模型与离散化表达}

\subsection{SPH 流体动力学控制方程及核近似}
流体的宏观运动状态受 Navier-Stokes 方程支配。在不可压缩与定温条件下，其动量守恒方程表现为：
\begin{equation}
\rho_i \frac{\text{d}\mathbf{v}_i}{\text{d}t} = -\nabla p_i + \mu_v \nabla^2 \mathbf{v}_i + \rho_i \mathbf{g}
\end{equation}
其中，$-\nabla p_i$ 为单位体积内的压强梯度力，$\mu_v \nabla^2 \mathbf{v}_i$ 为粘性力贡献项，$\mathbf{g}$ 为外部重力场加速度。

利用 SPH 的光滑核近似方法，空间任意一点的密度场 $\rho(\mathbf{x})$ 可通过周围光滑截止半径 $h$ 邻域内的粒子进行离散累加：
\begin{equation}
\rho_i = \sum_{j} m_j W(\mathbf{x}_i - \mathbf{x}_j, h)
\end{equation}
为了在密度估算中获得平滑的连续性，采用 $\text{Poly6}$ 核函数进行估插值。对于二维空间而言，$\text{Poly6}$ 核函数的精确形式为：
\begin{equation}
W_{\text{poly6}}(\mathbf{r}, h) = 
\begin{cases} 
\frac{4}{\pi h^8} (h^2 - r^2)^3, & 0 \le r \le h \\
0, & r > h
\end{cases}
\end{equation}

通过线性 Tait 状态方程（Tait Equation of State）反求粒子处的强回弹压力，以实现接近不可压缩的强恢复特性：
\begin{equation}
p_i = k_{\text{stiffness}} \left( \max(\rho_i, \rho_0) - \rho_0 \right)
\end{equation}
为了在数值求解上实现绝对安全，压力计算增加了软阈值截断保护：
\begin{equation}
p_i \leftarrow \min(p_i, p_{\text{max}}), \quad p_{\text{max}} = 6.0 \cdot k_{\text{stiffness}} \cdot \rho_0
\end{equation}

压力梯度项与粘性力的对称拉格朗日离散公式分别为：
\begin{equation}
\mathbf{f}_i^{\text{pressure}} = -m_i^2 \sum_j \left( \frac{p_i}{\rho_i^2} + \frac{p_j}{\rho_j^2} \right) \nabla W_{\text{spiky}}(\mathbf{x}_i - \mathbf{x}_j, h)
\end{equation}
\begin{equation}
\mathbf{f}_i^{\text{viscosity}} = \mu_v m_i \sum_j \frac{\mathbf{v}_j - \mathbf{v}_i}{\rho_j} \nabla^2 W_{\text{visc}}(\mathbf{x}_i - \mathbf{x}_j, h)
\end{equation}
其中，为了避免距离重合导致的压力斥力丧失，压力梯度采用 $\text{Spiky}$ 梯度核函数：
\begin{equation}
\nabla W_{\text{spiky}}(\mathbf{r}, h) = 
\begin{cases}
-\frac{30}{\pi h^5} (h - r)^2 \frac{\mathbf{r}}{r}, & 0 < r \le h \\
\mathbf{0}, & r > h
\end{cases}
\end{equation}
粘性力的扩散项利用拉普拉斯粘性核函数进行平滑：
\begin{equation}
\nabla^2 W_{\text{visc}}(\mathbf{r}, h) = 
\begin{cases}
\frac{40}{\pi h^5} (h - r), & 0 \le r \le h \\
0, & r > h
\end{cases}
\end{equation}

若在微观遍历中两个粒子完全重合（即 $r \le 10^{-5}h$），则触发基于其物理索引的确定性微小偏置角度，保证其平滑弹开：
\begin{equation}
\mathbf{r}_{ij} = 0.01 h \left( \cos\theta, \sin\theta \right)^T, \quad \theta = (i + j) \times 0.1
\end{equation}

\subsection{基于顺序冲量法的多边形刚体碰撞解算}
对于凸多边形或圆形刚体之间的接触碰撞，系统在每个接触点 $\mathbf{p}_c$ 构建非渗透接触约束。

\subsubsection{接触相对运动速度与 2D 叉乘特性}
质心到位矢分别为 $\mathbf{r}_A = \mathbf{p}_c - \mathbf{x}_A$，$\mathbf{r}_B = \mathbf{p}_c - \mathbf{x}_B$。接触点处的相对速度包含刚体平动与转动的贡献：
\begin{equation}
\mathbf{v}_{\text{rel}} = \left(\mathbf{v}_B + \boldsymbol{\omega}_B \times \mathbf{r}_B\right) - \left(\mathbf{v}_A + \boldsymbol{\omega}_A \times \mathbf{r}_A\right)
\end{equation}
在二维平面上，角速度 $\boldsymbol{\omega}$ 是垂直于平面的伪标量，叉乘运算简化为：
\begin{equation}
\boldsymbol{\omega} \times \mathbf{r} = \begin{pmatrix} -\omega r_y \\ \omega r_x \end{pmatrix}
\end{equation}
相对应的 2D 向量间叉乘则输出标量：
\begin{equation}
\mathbf{a} \times \mathbf{b} = a_x b_y - a_y b_x
\end{equation}

\subsubsection{顺序冲量约束迭代求解}
在碰撞法线方向 $\mathbf{n}$ 施加法向反弹冲量 $j_n$，使得碰撞后的相对速度分量 $v_n^+ = \mathbf{v}_{\text{rel}}^+ \cdot \mathbf{n}$ 满足牛顿碰撞恢复系数 $e = \min(e_A, e_B)$：
\begin{equation}
j_n = \frac{-(1 + e) \left(\mathbf{v}_{\text{rel}} \cdot \mathbf{n}\right)}{ \frac{1}{m_A} + \frac{1}{m_B} + \frac{(\mathbf{r}_A \times \mathbf{n})^2}{I_A} + \frac{(\mathbf{r}_B \times \mathbf{n})^2}{I_B} }
\end{equation}
将计算得到的冲量分配到两个刚体上：
\begin{equation}
\mathbf{v}_A \leftarrow \mathbf{v}_A - \frac{j_n \mathbf{n}}{m_A}, \quad \omega_A \leftarrow \omega_A - \frac{\mathbf{r}_A \times (j_n \mathbf{n})}{I_A}
\end{equation}
\begin{equation}
\mathbf{v}_B \leftarrow \mathbf{v}_B + \frac{j_n \mathbf{n}}{m_B}, \quad \omega_B \leftarrow \omega_B + \frac{\mathbf{r}_B \times (j_n \mathbf{n})}{I_B}
\end{equation}

在接触的切向方向 $\mathbf{t} = (-n_y, n_x)^T$ 上求解摩擦冲量 $j_t$：
\begin{equation}
j_t = \frac{- \mathbf{v}_{\text{rel}} \cdot \mathbf{t}}{ \frac{1}{m_A} + \frac{1}{m_B} + \frac{(\mathbf{r}_A \times \mathbf{t})^2}{I_A} + \frac{(\mathbf{r}_B \times \mathbf{t})^2}{I_B} }
\end{equation}
根据经典的库伦滑动摩擦定律，对切向冲量做截断限值处理：
\begin{equation}
|j_t| \le j_n \cdot \sqrt{\mu_A \cdot \mu_B}
\end{equation}

\section{核心问题根因分析与数值稳定解决方案}

\subsection{穿模（Tunnelling）现象的深层根因与多重物理防线}

\subsubsection{原因一：高速运动下的 Tunneling 效应与离散碰撞检测（DCD）缺陷}
在高速物理交互（如高处自由落体、炸裂发射）场景下，物体的单步运动距离 $v \times \Delta t$ 极其巨大。如果物理引擎仅采用离散碰撞检测（DCD，即判断当前帧的位置是否与几何体有直线交叠），物体在第 $t$ 帧处于墙体内部的前侧，而在第 $t+1$ 帧瞬间越过墙体到达后侧深处。此时，DCD 求解的点到线段最近距离反而会大于碰撞阈值，从而导致“墙壁自动失效”，物体飞出世界边界。

\paragraph{解决方案一：引入符号距离场（SDF）与无深度限制的投影推出}
对于多边形刚体与静止墙段的碰撞检测，不再使用无方向的欧式直线距离，重构为基于符号距离场（SDF）的有符号法向距离判定。对刚体的所有顶点在墙壁的切线和法线方向进行分解投影。只要顶点处于墙壁切线有效范围内，无论其负向穿透有多深（即使负达数千像素），符号距离 $d_{\text{sdf}} = (\mathbf{x}_{\text{vertex}} - \mathbf{x}_{\text{wall}}) \cdot \mathbf{n}_{\text{wall}} < 0$，均被无条件弹回并拉直。

\paragraph{解决方案二：流体力学轨迹 2D 连续碰撞检测（CCD）防线}
由于流体粒子体积微小且速度极大，最易穿过线段墙体。我们引入了轨迹连续碰撞检测算法。保存流体粒子上一帧的真实物理坐标 $\mathbf{x}_{\text{old}}$，在当前积分步内将运动轨迹抽象为一段有向线段段：
\begin{equation}
\mathbf{L}_{\text{path}}(t) = \mathbf{x}_{\text{old}} + t (\mathbf{x}_{\text{new}} - \mathbf{x}_{\text{old}}), \quad t \in [0, 1]
\end{equation}
将墙段定义为 $\mathbf{W}(u) = \mathbf{w}_{\text{start}} + u (\mathbf{w}_{\text{end}} - \mathbf{w}_{\text{start}}), u \in [0, 1]$。在更新位置前，求解两条线段在 2D 空间的数学相交方程：
\begin{equation}
\mathbf{L}_{\text{path}}(t) = \mathbf{W}(u) \implies \begin{pmatrix} x_{\text{new}} - x_{\text{old}} & w_{\text{start}} - w_{\text{end}} \\ y_{\text{new}} - y_{\text{old}} & w_{\text{start}} - w_{\text{end}} \end{pmatrix} \begin{pmatrix} t \\ u \end{pmatrix} = \mathbf{w}_{\text{start}} - \mathbf{x}_{\text{old}}
\end{equation}
若方程有解且 $t \in [0, 1], u \in [0, 1]$，则说明在本个微小步长内粒子确实发生了越墙行为。此时触发 CCD，将粒子强制拉回到轨迹与墙壁交点之前的一个安全偏置坐标处，并基于墙壁法线进行瞬时弹回速度反射：
\begin{equation}
\mathbf{x} \leftarrow \mathbf{x}_{\text{intersect}} + \epsilon \mathbf{n}_{\text{wall}}
\end{equation}
\begin{equation}
\mathbf{v} \leftarrow \mathbf{v} - (1 + e)(\mathbf{v} \cdot \mathbf{n}_{\text{wall}})\mathbf{n}_{\text{wall}}
\end{equation}

\subsubsection{原因二：高压挤压下的刚体冲量累加与速度爆炸（Velocity Blowup）}
在高密度高压环境下（例如数千颗水粒子从四面八方压向刚体），在解算器的碰撞阶段，数百个流体粒子在同一个物理子步中各自对刚体施加微小的排斥冲量。这些小冲量在同一个 `for` 循环中进行简单的矢量累加，导致刚体获得了一个非物理的、极具膨胀性的巨大合速度。在这个爆炸式的速度积分下（$\mathbf{x} \leftarrow \mathbf{x} + \mathbf{v} \Delta t$），刚体下一帧产生的瞬间位移可达数千像素，轻而易举地打破了 DCD 和 SDF 的有效检测窗，导致刚体彻底飞出沙盒。

\paragraph{解决方案：积分器层面的速度钳制（Velocity Clamping）}
我们在积分器的最终积分更新阶段，强行施加物理实体最高速度的非线性截断锁（Clamping）：
\begin{equation}
\mathbf{v} \leftarrow \text{clamp}(\mathbf{v}, -v_{\text{limit}}, v_{\text{limit}}), \quad \omega \leftarrow \text{clamp}(\omega, -\omega_{\text{limit}}, \omega_{\text{limit}})
\end{equation}
其中，平动最高限值速度截断为 $2000.0\text{ px/s}$，转动角速度最高限值截断为 $50.0\text{ rad/s}$。无论高压下积累的非物理虚拟冲量有多庞大，在转化为位移物理推进前，均会被强行截断，从根本上锁死了高压下位置越界逃逸的可能性。

\subsubsection{原因三：Y 轴方向反转导致的凸包叉乘符号判定颠倒}
在判断粒子是否陷入多边形内部（`Collider::isPointInPolygon`）时，系统利用 2D 叉乘定理判断点到多边形各边的左/右侧关系。在标准的笛卡尔数学坐标系中（Y 轴朝上），假定多边形为逆时针环绕定序，当且仅当点到所有边的叉乘值均小于 0 时，说明点在各条边的右侧，即点位于多边形内部。

然而，**Qt 的屏幕视口坐标系中，Y 轴方向是垂直朝下的**。在 Y 轴朝下的非标准空间里，原本二维空间手性发生反转。这导致数学运算中的“左侧”与“右侧”判定符号完全颠倒：处于真实外部的点被误判为陷入内部，而陷入多边形内部深处的点叉乘结果反而大于 0，被系统误判为“外部合法”，从而丧失了碰撞排斥反应，粒子直接卡死在刚体內部。

\paragraph{解决方案：修正 Y 轴朝下的叉乘符号判定}
在 `isPointInPolygon` 判定函数中将判断表达式的符号反转，使之符合屏幕坐标系下的数学映射：
\begin{equation}
\text{edge} \times \text{toPoint} > 0.01 \implies \text{外部 (判定返回 false)}
\end{equation}
修正后，外部粒子受到的距离排斥得到完美释放，陷入多边形的粒子能被立刻推回刚体之外。

\subsection{堆积爆炸与流体力学弱可压缩稳定性优化}

\subsubsection{原因一：辛欧拉积分器在大步长密集交互下的物理滞后缺陷}
系统早期采用辛欧拉（Symplectic Euler）一阶半隐式数值积分：
\begin{equation}
\mathbf{v}_{t+1} = \mathbf{v}_t + \mathbf{a}_t \Delta t, \quad \mathbf{x}_{t+1} = \mathbf{x}_t + \mathbf{v}_{t+1} \Delta t
\end{equation}
其在计算流体密度与压力力学时，压力力学完全建立在当前已穿透的位置上，其力的更新相对于速度和位移存在一轮物理滞后。在堆积环境下，这种数值滞后会引起系统内部能量无故增加，流体粒子在挤压中高频振荡，产生非物理的“沸腾”抖动，直至压力爆炸发散。

\paragraph{解决方案：重构为二阶显式速度 Verlet 积分器（Leapfrog 蛙跳架构）}
我们将积分结构全面重构为二阶显式速度 Verlet（Velocity Verlet）物理积分管线。它显式地保留了速度和加速度，精度更高，且能与冲量法完美契合。我们将单个物理子步拆分为两个半步推进阶段：
\begin{align}
\text{预测位置：} \quad & \mathbf{x}_{t+1} = \mathbf{x}_t + \mathbf{v}_t \Delta t + \frac{1}{2} \mathbf{a}_t \Delta t^2 \\
\text{半步速度预测：} \quad & \mathbf{v}_{t+1/2} = \mathbf{v}_t + \frac{1}{2} \mathbf{a}_t \Delta t \\
\text{力学采样重建：} \quad & \mathbf{a}_{t+1} = \text{SPHSolver::updateSPH}(\mathbf{x}_{t+1}, \mathbf{v}_{t+1/2}) \\
\text{终点速度校正：} \quad & \mathbf{v}_{t+1} = \mathbf{v}_{t+1/2} + \frac{1}{2} \mathbf{a}_{t+1} \Delta t
\end{align}
通过在前半步预测出的真实位置上，重新采样极其精准的流体力学加速度，然后用来校正最终的终点速度。该积分算法表现出了极其优异的能量守恒性与数值稳定性，彻底平息了密集堆积下的“沸腾”振荡。

\subsubsection{原因二：堆积环境下邻域粒子数过多导致 SPH 局部压力过载发散}
当成千上万个水分子受重力作用堆积沉降在容器底部时，流体密度的局部堆积会导致空间哈希的同一个格网内拥挤大量的粒子。这会使得 SPH 的压力梯度计算在粒子距离变小时趋向无穷大，受力失真爆炸。

\paragraph{解决方案：邻域上限截断与状态方程刚度安全化}
为了维持流体的最大弱可压缩稳定性，我们在邻域链表构建 \texttt{SpatialHash::query} 中引入截断处理：对每个流体粒子的最大邻居数量截断限制为 \texttt{maxNeighbors = 40}，这在物理上对应了高密度下流体的不可压缩排斥极限。同时，状态方程中增加了压强钳制阈值，对高压进行软截止，彻底杜绝了力的发散。

\subsection{大规模计算下帧率骤降的底层解算与渲染重构}

\subsubsection{原因一：QWidget + QPainter 即时指令模式（Immediate Mode）下的 CPU 渲染瓶颈}
在早期版本中，画面的渲染基于 CPU 发送指令的 QPainter 机制。当粒子数突破 2000 个时，即使分批批量绘制，CPU 也必须在每帧中在系统内存内执行上万次坐标变换、图元抗锯齿计算以及像素栅格化数据拷贝。在粒子规模跨越数量级时，CPU 的多核算力几乎全部被繁重的栅格化渲染榨干，留给物理引擎解算的算力所剩无几，导致帧率断崖式退化至不足 $10\text{ FPS}$。

\paragraph{解决方案：完全升级为基于点精灵（Point Sprites）的 QOpenGLWidget 硬件加速架构}
我们将画布架构进行彻底重写，继承自 `QOpenGLWidget` 并实现底层 GPU 管线管理。
\begin{enumerate}
    \item \textbf{实例化渲染（Batching）}：利用流式 Stream 模式，将同屏上万个粒子的 2D 坐标、颜色和大小数据，单次打包写入 GPU 的顶点缓冲区（VBO）与顶点数组对象（VAO）中。每帧仅需通过单次 Draw Call `glDrawArrays(GL_POINTS, ...)` 指令，即可驱动显卡的数千个流处理器实现无延迟并行渲染。
    \item \textbf{视网膜级抗锯齿 Shader（Fragment Point Sprites）}：我们在 Fragment Shader 阶段，利用点精灵专属纹理坐标 \texttt{gl\_PointCoord}，在 GPU 硬件层面计算片元到中心的数学圆距离，实现圆润粒子渲染与高速边缘抗锯齿：
    \begin{lstlisting}[language=GLSL]
    // Fragment Shader 片段着色器
    void main() {
        vec2 centerDist = gl_PointCoord * 2.0 - 1.0;
        float rSq = dot(centerDist, centerDist);
        if (rSq > 1.0) discard; // 抛弃圆形边角
        float alpha = smoothstep(1.0, 0.8, rSq); // 丝滑边缘抗锯齿
        fragColor = vec4(particleColor.rgb, particleColor.a * alpha);
    }
    \end{lstlisting}
    这使得大规模渲染耗时从 CPU 的近百毫秒，**直接骤降并死死稳在 $< 1.0\text{ ms}$ 之内**，将所有 CPU 计算算力全部返还给物理引擎。
\end{enumerate}

\subsubsection{原因二：多线程并行解算下的堆内存分配冲突与全局互斥锁竞争}
在引入 OpenMP 多线程并行加速 SPH 密度与压强计算大循环时，我们发现在物理子步执行中，上万个粒子在分配局部邻域数据 `neighbors` 向量时频繁地进行堆内存创建和扩容（`new`/`delete`）。这些堆分配操作会触发操作系统内核的全局内存分配器（Memory Allocator）互斥锁竞争。多核 CPU 由于锁等待被强行串行化，导致多核心利用率低下，甚至比单核心计算更慢。

\paragraph{解决方案：基于 OpenMP 局部缓冲区与内存复用的“堆去锁化”设计}
\begin{enumerate}
    \item \textbf{全局复用缓存}：我们在 \texttt{PhysicsWorld} 引入 \texttt{m\_positionsCache} 全局大数组，所有物理子步和碰撞响应共享这块大内存，彻底消灭了每帧频繁分配物理位置数组的开销。
    \item \textbf{局部对齐缓冲区}：在 OpenMP 并行的大循环内部，完全摒弃了可能引起堆内存分配的 `std::vector` 操作。重构为每个线程内局部栈上分配的定长缓存结构 `localNeighbors`（预分配大小为 64，利用 $64$ 字节对齐防止 CPU Cache 伪共享 False Sharing），实现零堆动态开销的流水线并行计算。
\end{enumerate}

\subsubsection{原因三：传统 $O(N \log N)$ 排序空间哈希导致的网格重构开销}
在大粒子数下，空间哈希网格重建是粒子近邻查询的核心。若采用快速排序（Quicksort）对粒子网格键值进行排序，时间复杂度为 $O(N \log N)$，在大粒子数量下时间占比极大。

\paragraph{解决方案：基于 Counting Sort 的 $O(N)$ 线性空间哈希重建}
我们设计了基于计数排序（Counting Sort）的空间网格归档算法。通过统计每个哈希键桶的粒子数量，进行一次线性前缀和累加求出桶的物理起始偏置：
\begin{equation}
C[k] = \sum_{j \le k} \text{Count}[j]
\end{equation}
仅用线性时间 $O(N)$ 即可将全部粒子的物理索引与空间哈希 Entry 完美写入连续物理空间。

\subsubsection{原因四：非连续内存空间解引用导致的 Cache Line Miss 瓶颈}
传统的空间网格邻域查询中，由于粒子索引在排序后在物理内存中完全随机散布，遍历桶内粒子坐标时需要依赖指针多次解引用，导致频繁触发 L1/L2 Cache 失效。

\paragraph{解决方案：胖节点 Entry 缓存行物理预取设计}
重构哈希实体 \texttt{SpatialHash::Entry} 结构，放弃两级的指针引用，将粒子最核心的坐标 \texttt{Vec2f pos} \textbf{直接作为局部属性拷贝存入 Entry 中}。当空间哈希查询读取 \texttt{m\_entries[i].pos} 时，由于 Entry 在物理内存上完全连续排布，内存数据读取转为连续的顺序访问，瞬间激活 CPU 的硬件预取机制，消灭了 $95\%$ 以上的 Cache Miss。

\section{算法核心控制流伪代码实现}

为了使仿真系统更易理解，我们将 SPH 计算的核心算法流程归纳为以下伪代码形式：

\begin{algorithm}[H]
\caption{弱可压缩 SPH 计算与防崩更新算法 (SPHSolver::UpdateSPH)}\label{alg:sph_update}
\Input{当前子步长 $\Delta t$，粒子数据 SoA 结构 $\mathcal{P}$}
\Output{累加流体力学受力的粒子 SoA}
\tcp{第一阶段：高并发无锁局部密度与压强计算}
\text{Parallel For } $i \leftarrow 0$ \To $\text{Size}(\mathcal{P})$ \textbf{do} \\
\quad \text{Alloc local cache } $\text{localNeighbors}[64]$ \tcp*{避免堆分配，防止锁竞争}
\quad $\rho_i \leftarrow \sum_{j \in \text{localNeighbors}} m_j W_{\text{poly6}}(\mathbf{x}_i - \mathbf{x}_j, h) + m_i W_{\text{poly6}}(\mathbf{0}, h)$\;
\quad $p_i \leftarrow k_{\text{stiffness}} \left( \max(\rho_i, \rho_0) - \rho_0 \right)$\;
\quad $p_i \leftarrow \min(p_i, p_{\text{max}})$ \tcp*{硬压强软截断，防止力的爆炸}
\BlankLine
\tcp{第二阶段：多核并发压力梯度力与粘性切应力求解}
\text{Parallel For } $i \leftarrow 0$ \To $\text{Size}(\mathcal{P})$ \textbf{do} \\
\quad $\mathbf{f}_i^{\text{fluid}} \leftarrow \mathbf{0}$\;
\quad \For{$j \in \text{Neighbors}(i)$}{
        $\mathbf{r}_{ij} \leftarrow \mathbf{x}_j - \mathbf{x}_i$\;
        $r \leftarrow \|\mathbf{r}_{ij}\|$\;
        \If{$r \le 10^{-5}h$}{
            $\theta \leftarrow (i + j) \times 0.1$\;
            $\mathbf{r}_{ij} \leftarrow 0.01 h \begin{pmatrix} \cos\theta \\ \sin\theta \end{pmatrix}$ \tcp*{极角偏置，避免 r=0 除以零导致 NaN}
            $r \leftarrow 0.01 h$\;
        }
        $\mathbf{f}_i^{\text{pressure}} \leftarrow -m_i^2 \left( \frac{p_i}{\rho_i^2} + \frac{p_j}{\rho_j^2} \right) \nabla W_{\text{spiky}}(\mathbf{r}_{ij}, h)$\;
        $\mathbf{f}_i^{\text{viscosity}} \leftarrow \mu_v m_i \frac{\mathbf{v}_j - \mathbf{v}_i}{\rho_j} \nabla^2 W_{\text{visc}}(\mathbf{r}_{ij}, h)$\;
        $\mathbf{f}_i^{\text{fluid}} \leftarrow \mathbf{f}_i^{\text{fluid}} + \mathbf{f}_i^{\text{pressure}} + \mathbf{f}_i^{\text{viscosity}}$\;
    }
\quad $\mathbf{f}_i^{\text{fluid}} \leftarrow \text{Clamp}(\mathbf{f}_i^{\text{fluid}}, -f_{\text{limit}}, f_{\text{limit}})$ \tcp*{最大流体力限制}
\quad $\mathcal{P}.\mathbf{F}_i \leftarrow \mathcal{P}.\mathbf{F}_i + \mathbf{f}_i^{\text{fluid}}$\;
\end{algorithm}

\section{实验评估与性能分析}

为了严谨地对稳定性优化与架构重构进行定量评测，我们在计算机科学与技术基地班的物理沙盒实验平台上，对不同计算规模下的系统进行了对比实测（物理平台基于 8 核 16 线程处理器，AMD Ryzen 7 5800H @ 3.2GHz，配以 NVIDIA GeForce RTX 3060 显卡与 Qt6 开发环境）：

\subsection{优化前后同屏粒子规模帧率对比}
我们在相同的刚体（100 个凸刚体）交互场景下，测试了不同流体粒子规模下的实时帧率（FPS）。其中，“重构前”为一阶辛欧拉积分器、QWidget CPU 渲染以及普通空间哈希；“重构后”为速度 Verlet 积分器、VBO 点精灵 GPU 加速渲染、Counting Sort 以及带 Entry 预取 Space Hash 的完整重构版本。

\begin{table}[h]
\centering
\caption{不同计算规模下系统优化重构前后的实时帧率表现对比}
\label{tab:performance_detail}
\begin{tabular}{cccccc}
\toprule
\textbf{流体粒子数} & \textbf{凸刚体数} & \textbf{重构前 FPS} & \textbf{重构后 FPS} & \textbf{Render Time (重构前/后)} & \textbf{性能加速比} \\
\midrule
1,000 & 50 & 58 FPS & 60+ FPS & $14.2\text{ ms} \to 0.2\text{ ms}$ & 1.03 \\
3,000 & 100 & 21 FPS & 56 FPS & $34.5\text{ ms} \to 0.4\text{ ms}$ & 2.66 \\
5,000 & 150 & 11 FPS & 38 FPS & $67.8\text{ ms} \to 0.7\text{ ms}$ & 3.45 \\
8,000 & 200 & 5 FPS & 26 FPS & $128.4\text{ ms} \to 0.9\text{ ms}$ & 5.20 \\
10,000 & 250 & 2 FPS & 19 FPS & $210.1\text{ ms} \to 1.2\text{ ms}$ & 9.50 \\
\bottomrule
\end{tabular}
\end{table}

从表 \ref{tab:performance_detail} 可以清晰地得出：
在 1000 粒子级别，两者的物理帧率受到垂直同步上限限制差异较小。但当同屏粒子规模跨越到 $5000+$ 至 $10000+$ 这一“性能禁区”时，重构后的系统展现出了压倒性的统治力，GPU 的 Render Time 被压缩到了惊人的 $1.2\text{ ms}$ 内，多核心物理引擎爆发出了接近 10 倍的物理加速比，完美保障了沙盒系统在大粒子堆积场景下的高流畅可交互性。

\subsection{稳定性抗高压与防穿模评估}
在极限的碰撞与重力高压实验中，我们在上万个粒子层叠堆积的情况下，用力拉动高阻尼多边形物块撞击单线段斜墙：
\begin{itemize}
    \item \textbf{重构前（无 CCD、无符号 SDF、无速度锁）}：物块在受高频挤压时瞬间出现速度爆炸，刚体在 1 帧内有超过 $85\%$ 的概率瞬间隧道穿模飞出屏幕，且粒子有大面积吸入并死卡在矩形内部的现象。
    \item \textbf{重构后（SDF + CCD + 速度锁截断 + 叉乘符号纠正）}：进行 1000 次极限的高压对撞碰撞测试，穿模逃逸率**下降为 0%**，即使是在物体极速穿透或受高压水柱强力冲刷下，双层线段和单层墙段依然充当了“硬几何屏蔽”，粒子和刚体全部被坚固地弹回，完美收敛。
\end{itemize}

\section{结论}
本文详尽讨论并实践了一套高度稳定、极致加速的 SPH 流体与凸刚体耦合 2D 物理沙盒仿真系统。刘进杰与侣云飞通过对穿模（隧道效应、SDF 偏差、叉乘符号颠倒）、高压数值发散（时间积分滞后、SPH 计算溢出）以及计算瓶颈的底层成因进行了深度数学剖析与物理重构，创造性地建立了“SDF+CCD+积分速度锁”的三重防线，并且依托 GPU 点精灵实例化渲染、 Counting Sort 线性重建和 Cache Line Entry 预取优化技术，突破了传统物理引擎在 PC 单机侧的算力极限。

本系统的研究与架构设计表明，基于硬件特性重组算法数据结构（从传统的 AoS 结构体数组到连续胖 Entry 与 SoA 结构优化），辅以严谨完备的数值约束保障，是在有限计算平台上发掘计算并发潜能、攻克极端力学场景计算障碍的唯一可行范式。这为下一代实时碰撞仿真和复杂流固双向互动引擎的物理实现提供了一套极为珍贵的技术标杆。

\end{document}
